%--------------------------------------
\section{IU30 Pantalla de Modificación de Perfil de Personal}

\subsection{Objetivo}
	Permitir al Director General actualizar la información de identificación, nombre completo, contacto y dirección de residencia de un colaborador previamente registrado en el sistema.

\subsection{Diseño}
	Esta pantalla \IUref{IU30}{Pantalla de Modificación de Perfil de Personal} (ver figura~\ref{IU30}) presenta un formulario extenso dividido en tres grupos lógicos:
	\begin{enumerate}
		\item **Datos de Identificación:** Muestra la CURP del empleado (campo deshabilitado no modificable por seguridad), y permite editar el RFC, Sexo y la Fecha de Nacimiento.
		\item **Nombre Completo:** Cajas de entrada de texto para modificar el primer nombre (obligatorio), segundo nombre, primer apellido (obligatorio) y segundo apellido.
		\item **Contacto y Dirección:** Campos para actualizar el Correo Electrónico (obligatorio), Teléfono, Calle, Número Exterior, Colonia, Código Postal, Municipio y Estado.
	\end{enumerate}

\IUfig[.55]{default}{IU30}{Pantalla de Modificación de Perfil de Personal.}

\subsection{Salidas}
	\begin{Citemize}
		\item Formulario precargado con los datos actuales del colaborador consultados desde la base de datos.
	\end{Citemize}

\subsection{Entradas}
	\begin{itemize}
		\item RFC (Opcional, 13 caracteres máximo).
		\item Sexo (Selección obligatoria).
		\item Fecha de Nacimiento (Calendario obligatorio).
		\item Primer Nombre y Primer Apellido (Texto obligatorio).
		\item Correo electrónico (Texto obligatorio en formato de email).
		\item Datos de domicilio (Calle, Número exterior, Colonia, Municipio, Estado y Código Postal de 5 dígitos).
		\item Teléfono (Texto opcional).
	\end{itemize}

\subsection{Comandos}
	\begin{itemize}
		\item \IUbutton{Cancelar}: Aborta la edición del colaborador y redirige al usuario de vuelta al Dashboard del Director.
		\item \IUbutton{Guardar Todos los Cambios}: Envía el formulario para su validación en el servidor y actualización correspondiente en las tablas `persona` y `personal`.
	\end{itemize}

\subsection{Mensajes}
	\begin{Citemize}
		\item {\bf MSG-09} (Operación exitosa): Se muestra tras persistir correctamente los datos modificados en la base de datos.
		\item {\bf MSG-04} (Error de validación): Se muestra si algún campo obligatorio se envía vacío o la sintaxis de correo o RFC es incorrecta.
	\end{Citemize}
%--------------------------------------
